\pagenumbering{arabic}
\section{INTRODUCCIÓN}\label{cap:introduccion}

Al comienzo de cada año escolar, las familias colombianas reciben del colegio una lista de
útiles que deben conseguir completa y en pocos días. El gasto es alto, porque el 35\,\% de
las familias destina entre \$300.000 y \$500.000 por niño \citep{noauthor_resultados_2025},
y la presión de tiempo también, porque el 70\,\% no planea la compra con anticipación
\citep{noauthor_resultados_2024}. A eso se suma el esfuerzo de interpretar cada renglón,
buscarlo en varias tiendas y comparar precios y marcas, una carga mental que lleva a decidir
peor \citep{sweller_cognitive_2011}. Cuando la compra se retrasa o queda incompleta, quien
lo paga es el estudiante, que empieza el año sin los materiales que el colegio pidió para
sus clases.\par

En Bogotá hay más de 1,17 millones de estudiantes matriculados
\citep{noauthor_estadisticas_nodate}, y este proyecto estudió el problema de cerca en el
Colegio Agustiniano Ciudad Salitre. Una encuesta a sus acudientes, con 60 respuestas
válidas, mostró que el 65\,\% dedica cuatro horas o más a la compra, que el 68,3\,\% compra
en tres o más lugares y que la queja más frecuente es no encontrar el producto o que esté
agotado (43,3\,\%). Las entrevistas con directivos, docentes y acudientes confirmaron esa
fricción.

Las alternativas que existen hoy no resuelven ese recorrido completo. Panamericana y Éxito
publican catálogos por categoría \citep{panamericana_cotizador_2025,exito_listas_2025},
Primavera permite armar un carrito a mano sobre su propio catálogo
\citep{primavera_armatulista_2026} y MercadoLibre vende kits fijos que arma cada vendedor
\citep{mercadolibre_kits_2025}. En todos los casos es el acudiente quien debe traducir la
lista del colegio en productos. Plataformas como TeacherLists, en Estados Unidos, convierten la lista del salón en un carrito de compra, pero ninguna
lee listas en papel o en foto ni está disponible en Colombia.

Frente a este panorama, el presente proyecto describe a Tukit como una plataforma que le permite al acudiente conseguir la lista de útiles completa.
Funciona como un sistema de recomendación multiobjetivo con un modelo B2B2C, en el que el
colegio aporta la lista oficial y la familia hace la compra. El acudiente elige el colegio y
el grado de cada hijo, o sube su propia lista en imagen, PDF o texto libre, e indica su
presupuesto y sus preferencias. A cambio recibe hasta cuatro kits completos (Ahorrador,
Premium, Fan y A tu Medida) y puede cambiar
cualquier producto antes de pagar en línea. Para llegar a esos kits, el sistema recorre tres etapas.\par

\begin{itemize}
  \item \textbf{Percepción.} Un modelo multimodal de inteligencia artificial lee la lista y
        la convierte en renglones.
  \item \textbf{Normalización.} Cada renglón se identifica como un útil, un texto escolar,
        una licencia digital u otro tipo, y se separa la cantidad que pide el colegio del
        empaque en que se vende el producto. Luego, cada útil se busca en un catálogo de
        conceptos canónicos con una búsqueda híbrida. Un canal léxico corrige la ortografía
        y exige que lo pedido esté nombrado, y un canal semántico compara significados para
        reconocer sinónimos. Si el sistema no tiene certeza suficiente, el renglón pasa a
        revisión humana.
  \item \textbf{Optimización.} Un algoritmo calcula el frente de Pareto de los kits
        posibles, es decir, aquellos en los que no se puede bajar el costo sin ceder en la
        marca o en las preferencias. Si las combinaciones son pocas (hasta 4.096), las revisa
        todas; si son más, usa NSGA-II \citep{deb_fast_2002}. De ese frente sale el kit de
        cada perfil.
\end{itemize}

Esta estrategia toma de un esquema RAG la indexación del catálogo y la recuperación de los
productos, pero en lugar de que un modelo de lenguaje redacte la respuesta, los kits los
construye el optimizador y el equipo de Tukit los revisa antes de publicarlos
(\emph{human-in-the-loop}). Todo corre sobre una arquitectura de seis servicios con cinco
bases de datos, elegida entre tres alternativas con el método ATAM
\citep{bass_software_2021}.\par

El trabajo siguió una metodología híbrida contenedora organizada en cinco fases, una por cada
objetivo específico. SCRUM \citep{schwaber_scrum_2020} organizó la construcción en ciclos de
dos semanas, Design Thinking \citep{brown_design_2008} guió la propuesta de valor y el modelo
de negocio, RUP \citep{kruchten_rup_2003} aportó la gestión de riesgos y la trazabilidad, y
CRISP-DM \citep{chapman_crisp-dm_2000} fue el marco de los componentes que dependen de datos.
Los estándares IEEE 830 \citep{ieee_830_1998}, 1016 \citep{ieee_1016_2009} y 1058
\citep{ieee_1058_1998} definieron la forma de los requerimientos, el diseño y el plan del
proyecto. Aunque el piloto se hizo en un colegio de Bogotá, el modelo puede llevarse a otros
colegios y ciudades, porque sumar un colegio consiste en cargar y revisar su lista, sin
cambiar el sistema.\par